\chapter{Introduction}

\section{Contest}
Until a few years ago, performance improvements depended mainly on increasing CPU clock speeds. Today, the approach has changed completely, shifting towards hardware specialization.

Modern computers, from servers to embedded SoCs, don't rely on a single processor anymore. Instead, they integrate different specific accelerators:

\begin{itemize}
    \item \textbf{CPU}: Optimized for sequential tasks and low latency.
    \item \textbf{GPU}: Available as integrated (iGPU) or discrete (dGPU), in a lot of cases both at the same time, designed for massive parallel computing.
    \item \textbf{NPU} (Neural Processing Unit): A recently introduced unit, specialized in speeding up matrix operations for AI.
\end{itemize}

This variety of hardware has created a gap with the software. While hardware offers different resources for different workloads, traditional software tends to be static. It usually assigns tasks to a single device (often just the CPU or GPU) and ignores the other available resources. 

The result is an inefficient use of the system. We could get more throughput with the same efficiency, or be more efficient at the same speed, so the current challenge is \textit{ orchestration}: we need a system that can dynamically map each task to the right accelerator, using the entire heterogeneous topology of the system.

\clearpage 

\section{Objectives}
The main goal of this thesis was to design and implement a scheduler capable of dynamically assigning computational tasks, specifically streams of matrix multiplications, to the different accelerators found in modern architectures.

The work started with an essential preliminary phase. I focused on integrating and studying different compute backends, such as CUDA, SYCL, OpenVINO. This step was useful to create a uniform abstraction layer over the hardware and to verify the technical feasibility of the system.

After this initial phase, the work moved on to the actual development of the scheduler. The core of the project focused on the asymmetric orchestration logic, designed to decide in real-time which resource to use for each specific workload.

Finally, the system was validated with experimental results on a complex hybrid architecture. This setup included an Intel multicore CPU (with integrated GPU and NPU) and a discrete NVIDIA GPU. The tests used to verify the operation were streams of identical matrices, streams of heterogeneous matrices, and splitting larger tasks across multiple accelerators.

\section{Thesis Structure}
The work is organized into the following chapters:

\begin{itemize}
    \item \textbf{Chapter 2}: Introduces the hardware technologies (CPU, GPU, NPU) and the software backends used (CUDA, SYCL, OpenVINO, OpenBLAS).

    \item \textbf{Chapter 3}: Describes the logical design, divided into two phases: the creation of the benchmarking environment and the subsequent design of the scheduler.    

    \item \textbf{Chapter 4}: Details the code implementation. It covers the wrappers for hardware abstraction, the essential data structures for the scheduler (ringbuffer and performance map), and the techniques used for task splitting.

    \item \textbf{Chapter 5}: Analyzes the experimental results, validating the effectiveness of the scheduler on different workloads.

    \item \textbf{Chapter 6}: Draws conclusions on the work done and suggests possible future developments.
\end{itemize}
